% Retention Strategy Taxonomy
% Maps retention mechanisms by control surface for §5.3
\begin{tikzpicture}[
  fam/.style={draw, rounded corners, thick, align=center, text width=26mm, minimum height=14mm, fill=black!3},
  method/.style={draw, rounded corners, align=center, text width=42mm, minimum height=14mm, fill=black!1, font=\footnotesize},
  note/.style={align=center, fill=white, inner sep=2pt, font=\footnotesize},
  arrow/.style={-{Latex[length=2.5mm]}, thick, shorten <=2pt, shorten >=2pt}
]

% 定义中心偏移量和列间距
\def\colOneX{1.5}
\def\colTwoX{6.5}

% Row 1: Protection
\node[fam] (F1) at (\colOneX,  0) {Protection\\\footnotesize constrain parameter\\\footnotesize movement};
\node[method] (M1a) at (\colTwoX, 0) {EWC, Online EWC\\\footnotesize Fisher-based importance weights, incremental update};

% Row 2: Replay
\node[fam] (F2) at (\colOneX,  -2.2) {Replay\\\footnotesize buffer samples\\\footnotesize absorb interference};
\node[method] (M2a) at (\colTwoX, -2.2) {GEM, A-GEM, Experience Replay, MIR\\\footnotesize gradient projection, reservoir sampling,\\interference-ranked selection};

% Row 3: Admission
\node[fam] (F3) at (\colOneX, -4.4) {Admission\\\footnotesize select which samples\\\footnotesize enter bounded store};
\node[method] (M3a) at (\colTwoX, -4.4) {iCaRL, GDumb, StreamFP\\\footnotesize exemplar anchoring, greedy curation,\\future-utility ranking};

% Row 4: Compression and Elasticity
\node[fam] (F4) at (\colOneX, -6.6) {Compression and Elasticity\\\footnotesize cheaper representation,\\\footnotesize budget contraction};
\node[method] (M4a) at (\colTwoX, -6.6) {REMIND, DER, Ferret\\\footnotesize feature-space replay, distilled matching,\\elastic budget and pipeline-aware};

% Arrows
\draw[arrow] (F1.east) -- (M1a.west);
\draw[arrow] (F2.east) -- (M2a.west);
\draw[arrow] (F3.east) -- (M3a.west);
\draw[arrow] (F4.east) -- (M4a.west);

% Annotation (居中放置于 F 和 M 的中间区域)
\node[note, text width=100mm] at (4.0,-8.2) {\textbf{Open gap:} unified retention-store contract for\\
admission, replay, compression, and demotion.};

\end{tikzpicture}